\section{Exploitation Mechanics}
\label{sec:exploitation}

\subsection{Injecting \texttt{LD\_PRELOAD} from the Container}
\label{subsec:injection}

The mechanics of environment variable injection deserve careful examination: In the OCI container model, environment variables are defined in two places: the container image's configuration (the \texttt{Env} field in the image manifest) \cite{oci_image_spec} and the runtime configuration provided at container launch (via \texttt{docker run -e} or Kubernetes \texttt{env} fields in a Pod specification \cite{k8s_security_docs}).
 
Both mechanisms ultimately result in the environment variables being written into the container's OCI \texttt{config.json} \cite{oci_runtime_spec}. The vulnerable NVIDIA toolkit versions propagate this environment to the hook process through the parent process chain without filtering \cite{nvidia_advisory_2025}. The following diagram illustrates the inheritance path:

\begin{lstlisting}[language=bash, caption={Environment propagation chain from container to privileged hook.}]
Container Runtime (dockerd / containerd)
    |
    |-- reads config.json (contains LD_PRELOAD set by attacker)
    |
    v
runc (OCI runtime) -> forks hook process
    |
    |-- inherits environment from runtime process
    |   (including LD_PRELOAD from container config)
    |
    v
nvidia-ctk hook [UID=0, running on HOST]
    |
    v
ld.so reads LD_PRELOAD -> loads evil.so -> constructor() executes
\end{lstlisting}

The key insight is that \texttt{runc} (or the equivalent runtime) spawns the hook process via \texttt{fork()} and \texttt{execve()}, but the \texttt{envp} argument passed to \texttt{execve()} (which should contain a clean, controlled environment) instead contains the attacker-influenced container environment.

\subsection{Privileged Hook Execution on the Host}
\label{subsec:hook-exec}

A critical aspect of the vulnerability is the privilege level at which the hook executes. OCI prestart hooks run in the host's mount namespace and with the full privilege set of the container runtime daemon, which is typically root \cite{oci_runtime_spec}. This is architecturally necessary: the hook must mount GPU device files into the container's namespace (an operation that requires \texttt{CAP\_SYS\_ADMIN} and the ability to interact with the host's device filesystem) \cite{linux_capabilities_manpage}.
 
What makes this particularly dangerous is that the hook runs \textit{before} the container's security policies (seccomp profiles, AppArmor policies, dropped capabilities) have been applied to the container process. The hook itself is not subject to the container's security configuration (it is a host process, not a container process \cite{nvidia_ctk_docs}).
 
Furthermore, in most production deployments, the NVIDIA Container Toolkit's hook binary is not subject to any additional Mandatory Access Control policy \cite{nvidia_advisory_2025}. It runs as a plain root process with full Linux capabilities, making it an extremely high-value target. Once the attacker's library is loaded into this process, there are effectively no further kernel-level restrictions preventing arbitrary host access.

\subsection{Escalation to Host Root}
\label{subsec:escalation}

The privilege escalation achieved through this vulnerability is immediate and unconditional. The constructor function in the malicious shared library executes in the security context of the \texttt{nvidia-ctk} hook process \cite{nvd_cve_2025_23266}, which means:
 
\begin{itemize}
  \item \textbf{UID 0 (root)} on the host system.
  \item \textbf{Full Linux capability set}, including \texttt{CAP\_SYS\_ADMIN}, \texttt{CAP\_DAC\_OVERRIDE}, \texttt{CAP\_CHOWN}, and all others \cite{linux_capabilities_manpage}.
  \item \textbf{Host mount namespace access}, allowing reading and writing of any file on the host filesystem \cite{linux_namespaces_manpage}.
  \item \textbf{Host network namespace access} (depending on how the runtime is configured), potentially enabling network-level attacks against other hosts.
\end{itemize}
 
Practical payloads an attacker could deploy include: writing an SSH public key to \texttt{/root/.ssh/authorized\_keys} for persistent access; creating a SUID copy of \texttt{/bin/bash} in a world-writable directory, modifying \texttt{/etc/passwd} or \texttt{/etc/sudoers} to add an unprivileged account with root access, installing a kernel rootkit or exfiltrating sensitive host secrets (API keys, TLS certificates, other container data) \cite{wiz_nvidia_2024}.

\subsection{Realistic Attack Scenario: Cloud Multi-Tenant AI}
\label{subsec:scenario}

To contextualise the real-world severity of NVIDIAScape, consider the following realistic attack scenario in a cloud AI infrastructure:
 
\paragraph{Setting.} A cloud provider offers GPU-accelerated compute instances for AI model training and inference. Multiple customers share the same physical host, each running their workloads in isolated containers. The host runs the NVIDIA Container Toolkit (vulnerable version) to provide GPU access to all tenants.
 
\paragraph{Attacker.} A malicious customer (Tenant A) provisions a standard GPU container instance (a service available to any paying customer with no special vetting). They have full control over their container's environment variables and can mount volumes.
 
\paragraph{Attack sequence.}
\begin{enumerate}
  \item Tenant A writes a malicious shared library (\texttt{evil.so}) implementing a reverse shell or a persistent backdoor into a volume accessible from the host.
  \item Tenant A starts their container with \texttt{LD\_PRELOAD} pointing to \texttt{evil.so}.
  \item The NVIDIA hook executes on the host, loads \texttt{evil.so}, and the payload runs as root.
  \item The attacker now has root on the host. They can: read the memory, disk, and GPU state of all other tenants on the same host, exfiltrate proprietary model weights from Tenant B's running container, extract API keys and credentials from environment variables of other containers, pivot to the cloud provider's internal network via the host's network interfaces.
\end{enumerate}
 
\paragraph{Impact.} The breach extends far beyond the attacker's own container. All co-located tenants are compromised, the cloud provider's infrastructure is at risk, and the attacker has achieved lateral movement potential into the provider's broader internal network. \cite{nvd_cve_2025_23266, wiz_nvidia_2024}.
 Remember that all these problems started from a standard, unprivileged GPU container.